
% SGA 3，报告 VII_B；英文重建本第一轮的中文译本。
% 范围：依据文献局部页 58--59，第 2.5.3 节全文，止于 2.6 前。
% 合并阅读本页 571--572；重版印刷页 560--561；原版边码 524；编者注 103--106。
% 依据：Polo--Gille Exp7B-13oct24.pdf。
% SHA-256 D303690E0195654BBE7C3DAA5D89EA23E24CD709A3575217234AC19AE8836694。
% 参照了 Jacob Reinhold 的 CC BY 4.0 Markdown，作为对照来源。本段无图。
\subsubsection*{2.5.3}
\phantomsection
\label{sga3:VIIB:2.5.3}
设 \(k\) 是域。\footnote{\label{sga3:VIIB:nde:103}编号 \SGARefLink{sga3:VIIB:2.5.3}{2.5.3} 及 \SGARefLink{sga3:VIIB:2.5.3.A}{2.5.3.A}--\SGARefLink{sga3:VIIB:2.5.3.C}{2.5.3.C} 由编者添加。}
\medskip
\noindent\textbf{定义 2.5.3.A。}
\phantomsection
\label{sga3:VIIB:2.5.3.A}
若交换 \(k\)-群概形同构于
\[
  \operatorname{Spec}\mathbf H(G),
\]
其中 \(G\) 是无穷小交换形式 \(k\)-群，则称其为\emph{幺幂}群。\footnote{%
\label{sga3:VIIB:nde:104}这与幺幂 \(k\)-群概形的“通常”概念一致；参见\SGARefLink{sga3:VIIB:2}{第 2 节}末 \SGARefLink{sga3:VIIB:2.8}{2.8} 的补充。}
另一方面，按照\SGARefLink{sga3:ix:definition:1:1}{报告 IX，定义 1.1}，\footnote{\label{sga3:VIIB:nde:105}原文说：“若交换 \(k\)-群概形同构于 \(\operatorname{Spec}\mathbf H(G)\)，其中 \(G\) 为平展交换形式 \(k\)-群，就称其为乘法型。”编者在此给出了取自\SGARefLink{sga3:ix:definition:1:1}{报告 IX，1.1} 的“通常”定义，并证明它与上述条件等价；另见 \([DG70]\)，§IV.1，\SGARefLink{sga3:VIIB:2.2}{定理 2.2}。}
若存在在 \(\operatorname{Spec}k\) 上忠实平坦且拟紧的概形 \(S\)，使得 \(H_S\) 是可对角化 \(S\)-群，就称 \(k\)-群概形 \(H\) 为\emph{乘法型}；即对某个“抽象”阿贝尔群 \(M\)，有
\[
  H_S\simeq D_S(M)
\]
。由 fpqc 下降，这蕴含 \(H\) 是仿射且交换的。此外，由于可用 \(S\) 的某个点的剩余域替换 \(S\)，\(H\) 为乘法型当且仅当存在 \(k\) 的扩张 \(K\)，使 \(H_K\) 为可对角化 \(K\)-群。
\medskip
\noindent\textbf{命题 2.5.3.B。}
\phantomsection
\label{sga3:VIIB:2.5.3.B}
{\itshape 设 \(T\) 是仿射交换 \(k\)-群概形，\(k_s\) 是 \(k\) 的可分闭包。
\begin{enumerate}
\item[\textup{(i)}] 概形 \(T\) 为乘法型当且仅当其对偶 \(D(T)\) 是平展交换形式 \(k\)-群。
\item[\textup{(ii)}] 因此，\(T\) 为乘法型当且仅当
\[
  T\otimes_k k_s
\]
可对角化。
\end{enumerate}}
\textit{证明。}
令 \(\mathcal A\) 为 \(D(T)\) 的仿射代数，\(\mathcal A_0\) 为其在单位元处的局部分量。则 \(D(T)\) 在 \(k\) 上平展当且仅当
\[
  \mathcal A_0=k.
\]
若 \(K\) 是 \(k\) 的扩张，则 \(\mathcal A_0\widehat\otimes_k K\) 是局部的，因为它是 Artin 局部环的逆极限，故与局部分量
\[
  \bigl(\mathcal A\widehat\otimes_k K\bigr)_0.
\]
一致。此外，由于 \(D(T)\) 的构造与基变换交换\textup{（参见\ \SGARefLink{sga3:VIIB:2.2.2}{2.2.2}）}，还有
\[
  \mathcal A\widehat\otimes_k K
  \simeq
  \mathcal A\bigl(D(T_K)\bigr).
\]
\noindent\begin{minipage}{\linewidth}
\phantomsection
\label{sga3:VIIB:2.5.3.B:base-change-source-note}
\textit{原文说明。}\enspace
法文重版印作 \(\mathcal A\widehat\otimes_k K\simeq D(T_K)\)。左边是代数，右边是形式群。因此英文本采用上面在对象类型上正确的预期读法，并在这里保留印本公式。
\end{minipage}
设 \(T\) 为乘法型。则存在 \(k\) 的扩张 \(K\)，使得
\[
  \mathcal O(T)\otimes_k K
\]
作为 Hopf \(K\)-代数同构于某个阿贝尔群 \(M\) 的群代数 \(KM\)。由此可得 \(\mathcal A\widehat\otimes_k K\) 同构于赋予积拓扑的代数 \(K^M\)。于是
\[
  K
  =
  \bigl(\mathcal A\widehat\otimes_k K\bigr)_0
  =
  \mathcal A_0\widehat\otimes_k K.
\]
这蕴含 \(\mathcal A_0=k\)，故 \(D(T)\) 是平展的。
反过来，设 \(D(T)\) 平展。则
\[
  D(T)\widehat\otimes_k k_s
  =
  D(T_{k_s})
\]
是与某个群 \(M\) 对应的常值 \(k_s\)-群。因此它的协变双代数是群代数 \(k_sM\)\textup{（参见\ \SGARefLink{sga3:VIIB:2.5.A}{注 2.5.A}）}，故
\[
  \mathcal O(T)\otimes_k k_s=k_sM.
\]
这证明了 \(T\) 为乘法型，且由扩张 \(k\to k_s\) 分裂。命题得证。
由 \SGARefLink{sga3:VIIB:2.2.2}{2.2.2}、\SGARefLink{sga3:VIIB:2.5.1:corollary}{2.5.1} 及 \SGARefLink{sga3:VIIB:2.5}{2.5} 得到：
\medskip
\noindent\textbf{推论 2.5.3.C。}
\phantomsection
\label{sga3:VIIB:2.5.3.C}
{\itshape 设 \(k\) 是域，\(G\) 是仿射交换 \(k\)-群概形。
\begin{enumerate}
\item[\textup{(i)}] 群 \(G\) 包含一个乘法型子群 \(G_m\)，使得 \(G/G_m\) 为幺幂群。
\item[\textup{(ii)}] 当 \(k\) 完美时，还有一个从 \(G\) 到 \(G_m\) 的典范收缩，因而 \(G\) 是一个幺幂群与一个乘法型群的积。
\item[\textup{\SGAInlineRefTarget{sga3:VIIB:2.5.3.C:item:iii}(iii)}]
\footnote{\label{sga3:VIIB:nde:106}第\textup{\SGARefLink{sga3:VIIB:2.5.3.C:item:iii}{(iii)}} 项由编者添加，是\SGARefLink{sga3:VIIB:2.5}{命题 2.5} 的推论。}
设 \(k_s\) 是 \(k\) 的可分闭包，置
\[
  \Gamma=\operatorname{Gal}(k_s/k).
\]
乘法型 \(k\)-群概形的范畴与连续 \(\Gamma\)-模的范畴反等价。
\end{enumerate}}
